\section{La arquitectura Transformer}
\label{sec:arquitectura_transformer}

Las secciones anteriores construyeron los conceptos fundamentales: las
limitaciones de los modelos clásicos, el problema del lenguaje, y el
mecanismo de atención como solución. Esta sección ensambla esos conceptos
en la arquitectura completa del Transformer, describiendo cada bloque en
el orden en que los datos los atraviesan.

%------------------------------------------------------------
\subsection{Visión general del pipeline}
\label{subsec:pipeline}

Antes de entrar en los detalles de cada componente, conviene tener una
imagen clara del flujo completo de información. Dada una secuencia de
tokens de entrada, el Transformer realiza los siguientes pasos en orden:

\begin{enumerate}
    \item Cada token se convierte en un vector denso mediante una capa de
    \textbf{embeddings}.
    \item Se suma un vector de \textbf{codificación posicional} para
    inyectar información de orden.
    \item La secuencia pasa por $N$ \textbf{bloques encoder} apilados,
    cada uno con atención multi-cabeza y una red feed-forward.
    \item (En modelos seq2seq) Una secuencia de salida parcial se procesa
    en paralelo por $N$ \textbf{bloques decoder}, que combinan
    auto-atención enmascarada con atención cruzada sobre el encoder.
    \item Una proyección lineal y un softmax producen una distribución
    de probabilidad sobre el vocabulario para el siguiente token.
\end{enumerate}

%------------------------------------------------------------
\subsection{Codificación posicional}
\label{subsec:positional_encoding}

El mecanismo de atención, tal como se describió en la
Sección~\ref{sec:mecanismo_atencion}, es invariante a permutaciones: si se
reordena la secuencia, los pesos de atención cambian pero el modelo no
tiene forma de saber que el orden cambió, porque trata todos los tokens
de forma simétrica. Esto contradice una de las propiedades más importantes
del lenguaje: el orden importa.

La solución propuesta por Vaswani et al.~\cite{vaswani2017attention} es
sumar a cada embedding un vector de \textit{codificación posicional} $PE_t$
que depende únicamente de la posición $t$ del token en la secuencia:

\begin{equation}
    \bm{x}'_t = \bm{e}_t + PE_t
\end{equation}

La codificación utiliza funciones sinusoidales de distintas frecuencias
para cada dimensión del vector:

\begin{align}
    PE_{(t,\, 2i)}   &= \sin\!\left(\frac{t}{10000^{2i/d_{\text{model}}}}\right) \\[4pt]
    PE_{(t,\, 2i+1)} &= \cos\!\left(\frac{t}{10000^{2i/d_{\text{model}}}}\right)
\end{align}

\noindent donde $t$ es la posición del token y $i$ es el índice de la
dimensión. Las dimensiones bajas oscilan con frecuencia alta (capturan
relaciones locales) y las altas con frecuencia baja (capturan relaciones
globales), de forma análoga a cómo un sistema numérico posicional codifica
valores de distinta magnitud en cada dígito.

Esta codificación tiene tres propiedades convenientes: es determinista
(no requiere aprendizaje adicional), generaliza a secuencias más largas
que las vistas en entrenamiento, y permite al modelo inferir distancias
relativas entre posiciones a través de diferencias entre vectores.

% ============================================================
% fig_transformer_completo.tex: Arquitectura Transformer completa
% Uso: % ============================================================
% fig_transformer_completo.tex — Arquitectura Transformer completa
% Uso: % ============================================================
% fig_transformer_completo.tex — Arquitectura Transformer completa
% Uso: \input{figures/fig_transformer_completo}
% ============================================================
\begin{figure}[H]
\centering
\begin{tikzpicture}[
    blk/.style={draw, rounded corners, minimum width=3.0cm,
                minimum height=0.62cm, font=\footnotesize, align=center},
    enc/.style={blk, fill=blue!10, draw=blue!40},
    dec/.style={blk, fill=green!10, draw=green!50!black},
    norm/.style={blk, fill=orange!10, draw=orange!50,
                 minimum height=0.5cm},
    emb/.style={blk, fill=purple!10, draw=purple!40},
    out/.style={blk, fill=red!8, draw=red!40},
    add/.style={draw, circle, fill=yellow!20, draw=yellow!60!black,
                minimum size=0.42cm, font=\scriptsize\bfseries},
    arrow/.style={-{Stealth}, thick},
    skip/.style={-{Stealth}, thick, gray!50, dashed},
    lbl/.style={font=\tiny, gray}
]

%═══════════════════════════════════════════════════════════
% ENCODER (izquierda)
%═══════════════════════════════════════════════════════════
\node[font=\small\bfseries, blue!60] at (-3.2, 10.8) {ENCODER};

% Entrada
\node[lbl] (src) at (-3.2, -0.4) {secuencia de entrada};
\node[emb] (src_emb) at (-3.2, 0.55) {Input Embedding};
\node[emb] (src_pe)  at (-3.2, 1.45) {Positional Encoding};

% Bloque encoder (representado como caja grande)
\draw[blue!30, rounded corners, thick, fill=blue!3]
    (-4.9, 2.1) rectangle (-1.5, 7.9);
\node[lbl] at (-4.55, 7.65) {$\times N$};

\node[enc]  (mha_e)  at (-3.2, 2.75) {Multi-Head\\Self-Attention};
\node[add]  (add1e)  at (-3.2, 3.75) {$+$};
\node[norm] (ln1e)   at (-3.2, 4.4)  {Add \& Norm};
\node[enc]  (ffn_e)  at (-3.2, 5.3)  {Feed-Forward};
\node[add]  (add2e)  at (-3.2, 6.3)  {$+$};
\node[norm] (ln2e)   at (-3.2, 6.95) {Add \& Norm};

% Salida encoder
\node[lbl] (enc_out) at (-3.2, 8.3) {$K, V$ al decoder};

% Flechas encoder
\draw[arrow] (src)     -- (src_emb);
\draw[arrow] (src_emb) -- (src_pe);
\draw[arrow] (src_pe)  -- (mha_e);
\draw[arrow] (mha_e)   -- (add1e);
\draw[arrow] (add1e)   -- (ln1e);
\draw[arrow] (ln1e)    -- (ffn_e);
\draw[arrow] (ffn_e)   -- (add2e);
\draw[arrow] (add2e)   -- (ln2e);
\draw[arrow] (ln2e)    -- (enc_out);

% Residuales encoder
\draw[skip] (src_pe.west) -- ++(-0.7,0) -- ++(0,3.75) -- (add1e.west);
\draw[skip] (ln1e.west)   -- ++(-0.5,0) -- ++(0,2.0)  -- (add2e.west);

%═══════════════════════════════════════════════════════════
% DECODER (derecha)
%═══════════════════════════════════════════════════════════
\node[font=\small\bfseries, green!60!black] at (3.2, 10.8) {DECODER};

% Entrada
\node[lbl] (tgt) at (3.2, -0.4) {secuencia de salida (desplazada)};
\node[emb] (tgt_emb) at (3.2, 0.55) {Output Embedding};
\node[emb] (tgt_pe)  at (3.2, 1.45) {Positional Encoding};

% Bloque decoder
\draw[green!30, rounded corners, thick, fill=green!3]
    (1.5, 2.1) rectangle (4.9, 9.1);
\node[lbl] at (4.55, 8.85) {$\times N$};

\node[dec]  (mha_d1) at (3.2, 2.75) {Masked Multi-Head\\Self-Attention};
\node[add]  (add1d)  at (3.2, 3.75) {$+$};
\node[norm] (ln1d)   at (3.2, 4.4)  {Add \& Norm};
\node[dec]  (mha_d2) at (3.2, 5.3)  {Multi-Head\\Cross-Attention};
\node[add]  (add2d)  at (3.2, 6.3)  {$+$};
\node[norm] (ln2d)   at (3.2, 6.95) {Add \& Norm};
\node[dec]  (ffn_d)  at (3.2, 7.85) {Feed-Forward};
\node[add]  (add3d)  at (3.2, 8.75) {$+$};
% (ln3 dentro del bloque, salida)

% Flechas decoder
\draw[arrow] (tgt)     -- (tgt_emb);
\draw[arrow] (tgt_emb) -- (tgt_pe);
\draw[arrow] (tgt_pe)  -- (mha_d1);
\draw[arrow] (mha_d1)  -- (add1d);
\draw[arrow] (add1d)   -- (ln1d);
\draw[arrow] (ln1d)    -- (mha_d2);
\draw[arrow] (mha_d2)  -- (add2d);
\draw[arrow] (add2d)   -- (ln2d);
\draw[arrow] (ln2d)    -- (ffn_d);
\draw[arrow] (ffn_d)   -- (add3d);

% Residuales decoder
\draw[skip] (tgt_pe.east) -- ++(0.7,0) -- ++(0,3.75) -- (add1d.east);
\draw[skip] (ln1d.east)   -- ++(0.5,0) -- ++(0,2.0)  -- (add2d.east);
\draw[skip] (ln2d.east)   -- ++(0.5,0) -- ++(0,2.0)  -- (add3d.east);

% ── Conexión encoder → decoder (cross-attention) ─────────────
\draw[arrow, red!60, very thick]
    (enc_out.east) -- ++(0.5,0) -- ++(0,-3.0) -- (mha_d2.west);
\node[font=\tiny, red!60, align=center] at (0, 6.8)
    {$K$, $V$ del\\encoder};

% ── Salida final ──────────────────────────────────────────────
\node[out]  (linear)  at (3.2, 9.8)  {Linear};
\node[out]  (softmax) at (3.2, 10.5) {Softmax};
\node[lbl]           at (3.2, 11.1)  {probabilidades de salida};

\draw[arrow] (add3d)  -- (linear);
\draw[arrow] (linear)  -- (softmax);

% ── Nota máscara ─────────────────────────────────────────────
\node[draw=gray!30, fill=gray!5, rounded corners, font=\tiny,
      inner sep=4pt, align=left] at (6.8, 2.75)
    {\textbf{Masked:} m\'ascara causal\\impide ver tokens futuros.\\Garantiza generaci\'on\\autoregresiva.};
\draw[-{Stealth}, gray!40, thin] (4.9, 2.75) -- (5.5, 2.75);

\end{tikzpicture}
\caption{Arquitectura completa del Transformer original~\cite{vaswani2017attention}.
El encoder (izquierda) procesa la secuencia de entrada en paralelo y produce
representaciones contextualizadas. El decoder (derecha) genera la secuencia
de salida token a token, utilizando atenci\'on enmascarada sobre su propia
salida parcial y atenci\'on cruzada (\textit{cross-attention}) sobre las
representaciones del encoder. Las flechas rojas indican el flujo de $K$ y $V$
del encoder al decoder.}
\label{fig:transformer_completo}
\end{figure}

% ============================================================
\begin{figure}[H]
\centering
\begin{tikzpicture}[
    blk/.style={draw, rounded corners, minimum width=3.0cm,
                minimum height=0.62cm, font=\footnotesize, align=center},
    enc/.style={blk, fill=blue!10, draw=blue!40},
    dec/.style={blk, fill=green!10, draw=green!50!black},
    norm/.style={blk, fill=orange!10, draw=orange!50,
                 minimum height=0.5cm},
    emb/.style={blk, fill=purple!10, draw=purple!40},
    out/.style={blk, fill=red!8, draw=red!40},
    add/.style={draw, circle, fill=yellow!20, draw=yellow!60!black,
                minimum size=0.42cm, font=\scriptsize\bfseries},
    arrow/.style={-{Stealth}, thick},
    skip/.style={-{Stealth}, thick, gray!50, dashed},
    lbl/.style={font=\tiny, gray}
]

%═══════════════════════════════════════════════════════════
% ENCODER (izquierda)
%═══════════════════════════════════════════════════════════
\node[font=\small\bfseries, blue!60] at (-3.2, 10.8) {ENCODER};

% Entrada
\node[lbl] (src) at (-3.2, -0.4) {secuencia de entrada};
\node[emb] (src_emb) at (-3.2, 0.55) {Input Embedding};
\node[emb] (src_pe)  at (-3.2, 1.45) {Positional Encoding};

% Bloque encoder (representado como caja grande)
\draw[blue!30, rounded corners, thick, fill=blue!3]
    (-4.9, 2.1) rectangle (-1.5, 7.9);
\node[lbl] at (-4.55, 7.65) {$\times N$};

\node[enc]  (mha_e)  at (-3.2, 2.75) {Multi-Head\\Self-Attention};
\node[add]  (add1e)  at (-3.2, 3.75) {$+$};
\node[norm] (ln1e)   at (-3.2, 4.4)  {Add \& Norm};
\node[enc]  (ffn_e)  at (-3.2, 5.3)  {Feed-Forward};
\node[add]  (add2e)  at (-3.2, 6.3)  {$+$};
\node[norm] (ln2e)   at (-3.2, 6.95) {Add \& Norm};

% Salida encoder
\node[lbl] (enc_out) at (-3.2, 8.3) {$K, V$ al decoder};

% Flechas encoder
\draw[arrow] (src)     -- (src_emb);
\draw[arrow] (src_emb) -- (src_pe);
\draw[arrow] (src_pe)  -- (mha_e);
\draw[arrow] (mha_e)   -- (add1e);
\draw[arrow] (add1e)   -- (ln1e);
\draw[arrow] (ln1e)    -- (ffn_e);
\draw[arrow] (ffn_e)   -- (add2e);
\draw[arrow] (add2e)   -- (ln2e);
\draw[arrow] (ln2e)    -- (enc_out);

% Residuales encoder
\draw[skip] (src_pe.west) -- ++(-0.7,0) -- ++(0,3.75) -- (add1e.west);
\draw[skip] (ln1e.west)   -- ++(-0.5,0) -- ++(0,2.0)  -- (add2e.west);

%═══════════════════════════════════════════════════════════
% DECODER (derecha)
%═══════════════════════════════════════════════════════════
\node[font=\small\bfseries, green!60!black] at (3.2, 10.8) {DECODER};

% Entrada
\node[lbl] (tgt) at (3.2, -0.4) {secuencia de salida (desplazada)};
\node[emb] (tgt_emb) at (3.2, 0.55) {Output Embedding};
\node[emb] (tgt_pe)  at (3.2, 1.45) {Positional Encoding};

% Bloque decoder
\draw[green!30, rounded corners, thick, fill=green!3]
    (1.5, 2.1) rectangle (4.9, 9.1);
\node[lbl] at (4.55, 8.85) {$\times N$};

\node[dec]  (mha_d1) at (3.2, 2.75) {Masked Multi-Head\\Self-Attention};
\node[add]  (add1d)  at (3.2, 3.75) {$+$};
\node[norm] (ln1d)   at (3.2, 4.4)  {Add \& Norm};
\node[dec]  (mha_d2) at (3.2, 5.3)  {Multi-Head\\Cross-Attention};
\node[add]  (add2d)  at (3.2, 6.3)  {$+$};
\node[norm] (ln2d)   at (3.2, 6.95) {Add \& Norm};
\node[dec]  (ffn_d)  at (3.2, 7.85) {Feed-Forward};
\node[add]  (add3d)  at (3.2, 8.75) {$+$};
% (ln3 dentro del bloque, salida)

% Flechas decoder
\draw[arrow] (tgt)     -- (tgt_emb);
\draw[arrow] (tgt_emb) -- (tgt_pe);
\draw[arrow] (tgt_pe)  -- (mha_d1);
\draw[arrow] (mha_d1)  -- (add1d);
\draw[arrow] (add1d)   -- (ln1d);
\draw[arrow] (ln1d)    -- (mha_d2);
\draw[arrow] (mha_d2)  -- (add2d);
\draw[arrow] (add2d)   -- (ln2d);
\draw[arrow] (ln2d)    -- (ffn_d);
\draw[arrow] (ffn_d)   -- (add3d);

% Residuales decoder
\draw[skip] (tgt_pe.east) -- ++(0.7,0) -- ++(0,3.75) -- (add1d.east);
\draw[skip] (ln1d.east)   -- ++(0.5,0) -- ++(0,2.0)  -- (add2d.east);
\draw[skip] (ln2d.east)   -- ++(0.5,0) -- ++(0,2.0)  -- (add3d.east);

% ── Conexión encoder → decoder (cross-attention) ─────────────
\draw[arrow, red!60, very thick]
    (enc_out.east) -- ++(0.5,0) -- ++(0,-3.0) -- (mha_d2.west);
\node[font=\tiny, red!60, align=center] at (0, 6.8)
    {$K$, $V$ del\\encoder};

% ── Salida final ──────────────────────────────────────────────
\node[out]  (linear)  at (3.2, 9.8)  {Linear};
\node[out]  (softmax) at (3.2, 10.5) {Softmax};
\node[lbl]           at (3.2, 11.1)  {probabilidades de salida};

\draw[arrow] (add3d)  -- (linear);
\draw[arrow] (linear)  -- (softmax);

% ── Nota máscara ─────────────────────────────────────────────
\node[draw=gray!30, fill=gray!5, rounded corners, font=\tiny,
      inner sep=4pt, align=left] at (6.8, 2.75)
    {\textbf{Masked:} m\'ascara causal\\impide ver tokens futuros.\\Garantiza generaci\'on\\autoregresiva.};
\draw[-{Stealth}, gray!40, thin] (4.9, 2.75) -- (5.5, 2.75);

\end{tikzpicture}
\caption{Arquitectura completa del Transformer original~\cite{vaswani2017attention}.
El encoder (izquierda) procesa la secuencia de entrada en paralelo y produce
representaciones contextualizadas. El decoder (derecha) genera la secuencia
de salida token a token, utilizando atenci\'on enmascarada sobre su propia
salida parcial y atenci\'on cruzada (\textit{cross-attention}) sobre las
representaciones del encoder. Las flechas rojas indican el flujo de $K$ y $V$
del encoder al decoder.}
\label{fig:transformer_completo}
\end{figure}

% ============================================================
\begin{figure}[H]
\centering
\begin{tikzpicture}[
    blk/.style={draw, rounded corners, minimum width=3.0cm,
                minimum height=0.62cm, font=\footnotesize, align=center},
    enc/.style={blk, fill=blue!10, draw=blue!40},
    dec/.style={blk, fill=green!10, draw=green!50!black},
    norm/.style={blk, fill=orange!10, draw=orange!50,
                 minimum height=0.5cm},
    emb/.style={blk, fill=purple!10, draw=purple!40},
    outblk/.style={blk, fill=red!8, draw=red!40},
    add/.style={draw, circle, fill=yellow!20, draw=yellow!60!black,
                minimum size=0.42cm, font=\scriptsize\bfseries},
    arrow/.style={-{Stealth}, thick},
    skip/.style={-{Stealth}, thick, gray!50, dashed},
    lbl/.style={font=\tiny, gray}
]

%═══════════════════════════════════════════════════════════
% ENCODER (izquierda)
%═══════════════════════════════════════════════════════════
\node[font=\small\bfseries, blue!60] at (-3.2, 10.8) {ENCODER};

% Entrada
\node[lbl] (src) at (-3.2, -0.4) {secuencia de entrada};
\node[emb] (src_emb) at (-3.2, 0.55) {Input Embedding};
\node[emb] (src_pe)  at (-3.2, 1.45) {Positional Encoding};

% Bloque encoder (representado como caja grande)
\draw[blue!30, rounded corners, thick, fill=blue!3]
    (-4.9, 2.1) rectangle (-1.5, 7.9);
\node[lbl] at (-4.55, 7.65) {$\times N$};

\node[enc]  (mha_e)  at (-3.2, 2.75) {Multi-Head\\Self-Attention};
\node[add]  (add1e)  at (-3.2, 3.75) {$+$};
\node[norm] (ln1e)   at (-3.2, 4.4)  {Add \& Norm};
\node[enc]  (ffn_e)  at (-3.2, 5.3)  {Feed-Forward};
\node[add]  (add2e)  at (-3.2, 6.3)  {$+$};
\node[norm] (ln2e)   at (-3.2, 6.95) {Add \& Norm};

% Salida encoder
\node[lbl] (enc_out) at (-3.2, 8.3) {$K, V$ al decoder};

% Flechas encoder
\draw[arrow] (src)     -- (src_emb);
\draw[arrow] (src_emb) -- (src_pe);
\draw[arrow] (src_pe)  -- (mha_e);
\draw[arrow] (mha_e)   -- (add1e);
\draw[arrow] (add1e)   -- (ln1e);
\draw[arrow] (ln1e)    -- (ffn_e);
\draw[arrow] (ffn_e)   -- (add2e);
\draw[arrow] (add2e)   -- (ln2e);
\draw[arrow] (ln2e)    -- (enc_out);

% Residuales encoder
\draw[skip] (src_pe.west) -- ++(-0.7,0) -- ++(0,3.75) -- (add1e.west);
\draw[skip] (ln1e.west)   -- ++(-0.5,0) -- ++(0,2.0)  -- (add2e.west);

%═══════════════════════════════════════════════════════════
% DECODER (derecha)
%═══════════════════════════════════════════════════════════
\node[font=\small\bfseries, green!60!black] at (3.2, 10.8) {DECODER};

% Entrada
\node[lbl] (tgt) at (3.2, -0.4) {secuencia de salida (desplazada)};
\node[emb] (tgt_emb) at (3.2, 0.55) {Output Embedding};
\node[emb] (tgt_pe)  at (3.2, 1.45) {Positional Encoding};

% Bloque decoder
\draw[green!30, rounded corners, thick, fill=green!3]
    (1.5, 2.1) rectangle (4.9, 9.1);
\node[lbl] at (4.55, 8.85) {$\times N$};

\node[dec]  (mha_d1) at (3.2, 2.75) {Masked Multi-Head\\Self-Attention};
\node[add]  (add1d)  at (3.2, 3.75) {$+$};
\node[norm] (ln1d)   at (3.2, 4.4)  {Add \& Norm};
\node[dec]  (mha_d2) at (3.2, 5.3)  {Multi-Head\\Cross-Attention};
\node[add]  (add2d)  at (3.2, 6.3)  {$+$};
\node[norm] (ln2d)   at (3.2, 6.95) {Add \& Norm};
\node[dec]  (ffn_d)  at (3.2, 7.85) {Feed-Forward};
\node[add]  (add3d)  at (3.2, 8.75) {$+$};
% (ln3 dentro del bloque, salida)

% Flechas decoder
\draw[arrow] (tgt)     -- (tgt_emb);
\draw[arrow] (tgt_emb) -- (tgt_pe);
\draw[arrow] (tgt_pe)  -- (mha_d1);
\draw[arrow] (mha_d1)  -- (add1d);
\draw[arrow] (add1d)   -- (ln1d);
\draw[arrow] (ln1d)    -- (mha_d2);
\draw[arrow] (mha_d2)  -- (add2d);
\draw[arrow] (add2d)   -- (ln2d);
\draw[arrow] (ln2d)    -- (ffn_d);
\draw[arrow] (ffn_d)   -- (add3d);

% Residuales decoder
\draw[skip] (tgt_pe.east) -- ++(0.7,0) -- ++(0,3.75) -- (add1d.east);
\draw[skip] (ln1d.east)   -- ++(0.5,0) -- ++(0,2.0)  -- (add2d.east);
\draw[skip] (ln2d.east)   -- ++(0.5,0) -- ++(0,2.0)  -- (add3d.east);

% ── Conexión encoder → decoder (cross-attention) ─────────────
\draw[arrow, red!60, very thick]
    (enc_out.east) -- ++(0.5,0) -- ++(0,-3.0) -- (mha_d2.west);
\node[font=\tiny, red!60, align=center] at (0, 6.8)
    {$K$, $V$ del\\encoder};

% ── Salida final ──────────────────────────────────────────────
\node[outblk]  (linear)  at (3.2, 9.8)  {Linear};
\node[outblk]  (softmax) at (3.2, 10.5) {Softmax};
\node[lbl]           at (3.2, 11.1)  {probabilidades de salida};

\draw[arrow] (add3d)  -- (linear);
\draw[arrow] (linear)  -- (softmax);

% ── Nota máscara ─────────────────────────────────────────────
\node[draw=gray!30, fill=gray!5, rounded corners, font=\tiny,
      inner sep=4pt, align=left] at (6.8, 2.75)
    {\textbf{Masked:} m\'ascara causal\\impide ver tokens futuros.\\Garantiza generaci\'on\\autoregresiva.};
\draw[-{Stealth}, gray!40, thin] (4.9, 2.75) -- (5.5, 2.75);

\end{tikzpicture}
\caption{Arquitectura completa del Transformer original~\cite{vaswani2017attention}.
El encoder (izquierda) procesa la secuencia de entrada en paralelo y produce
representaciones contextualizadas. El decoder (derecha) genera la secuencia
de salida token a token, utilizando atenci\'on enmascarada sobre su propia
salida parcial y atenci\'on cruzada (\textit{cross-attention}) sobre las
representaciones del encoder. Las flechas rojas indican el flujo de $K$ y $V$
del encoder al decoder.}
\label{fig:positional_encoding}
\end{figure}


%------------------------------------------------------------
\subsection{El bloque encoder}
\label{subsec:bloque_encoder}

El bloque encoder es la unidad fundamental de procesamiento. Recibe una
secuencia de vectores y devuelve una secuencia del mismo tamaño, donde
cada vector es una representación contextualizada de su token
correspondiente. Está compuesto por dos sub-capas:

\begin{enumerate}
    \item \textbf{Multi-Head Self-Attention:} cada token consulta a todos
    los demás para construir su representación contextualizada, como se
    describió en la Sección~\ref{sec:mecanismo_atencion}.

    \item \textbf{Feed-Forward Network (FFN):} una red de dos capas
    lineales con activación no lineal, aplicada de forma idéntica e
    independiente a cada posición:
    \begin{equation}
        \text{FFN}(\bm{x}) = \max(0,\, \bm{x} W_1 + \bm{b}_1)\, W_2 + \bm{b}_2
    \end{equation}
    La FFN actúa como una memoria de patrones por posición: mientras que
    la atención mezcla información entre tokens, la FFN transforma cada
    representación de forma local.
\end{enumerate}

Ambas sub-capas están envueltas en dos mecanismos de estabilización:

\begin{itemize}
    \item \textbf{Conexión residual (\textit{skip connection}):} la entrada
    de cada sub-capa se suma a su salida antes de la normalización:
    $\text{salida} = \text{sub-capa}(\bm{x}) + \bm{x}$. Esto permite que
    los gradientes fluyan directamente hacia las capas anteriores durante
    el entrenamiento, resolviendo el problema del gradiente desvanecido
    que afectaba a las RNNs.

    \item \textbf{Layer Normalization:} normaliza las activaciones de cada
    vector de forma independiente, estabilizando el entrenamiento y
    acelerando la convergencia~\cite{goodfellow2016deeplearning}.
\end{itemize}

% ============================================================
% fig_encoder_block.tex: Bloque Encoder detallado
% Uso: % ============================================================
% fig_encoder_block.tex: Bloque Encoder detallado
% Uso: % ============================================================
% fig_encoder_block.tex: Bloque Encoder detallado
% Uso: \input{figures/fig_encoder_block}
% ============================================================
\begin{figure}[H]
\centering
\begin{tikzpicture}[
    block/.style={draw, rounded corners, fill=blue!8, draw=blue!40,
                  minimum width=4.2cm, minimum height=0.7cm,
                  font=\small, align=center},
    norm/.style={draw, rounded corners, fill=orange!10, draw=orange!50,
                 minimum width=4.2cm, minimum height=0.55cm,
                 font=\small, align=center},
    add/.style={draw, circle, fill=yellow!15, draw=yellow!60!black,
                minimum size=0.5cm, font=\small\bfseries},
    arrow/.style={-{Stealth}, thick},
    skip/.style={-{Stealth}, thick, gray!60, dashed}
]

% ── Entrada ──────────────────────────────────────────────────
\node[font=\footnotesize, gray] (inp) at (0, 0) {$\bm{X}$ (embeddings + PE)};

% ── Sub-bloque 1: Multi-Head Self-Attention ──────────────────
\node[block] (mha)  at (0, 1.4) {Multi-Head\\Self-Attention};
\node[add]   (add1) at (0, 2.6) {$+$};
\node[norm]  (ln1)  at (0, 3.4) {Layer Norm};

% ── Sub-bloque 2: Feed-Forward ───────────────────────────────
\node[block] (ffn)  at (0, 4.6) {Feed-Forward\\Network};
\node[add]   (add2) at (0, 5.8) {$+$};
\node[norm]  (ln2)  at (0, 6.6) {Layer Norm};

% ── Salida ───────────────────────────────────────────────────
\node[font=\footnotesize, gray] (out) at (0, 7.6) {$\bm{Z}$ (representaciones contextualizadas)};

% ── Flechas principales ──────────────────────────────────────
\draw[arrow] (inp)  -- (mha);
\draw[arrow] (mha)  -- (add1);
\draw[arrow] (add1) -- (ln1);
\draw[arrow] (ln1)  -- (ffn);
\draw[arrow] (ffn)  -- (add2);
\draw[arrow] (add2) -- (ln2);
\draw[arrow] (ln2)  -- (out);

% ── Conexiones residuales ─────────────────────────────────────
% Residual 1: entrada → add1
\draw[skip] (inp.east) -- ++(1.6,0) -- ++(0,2.6) -- (add1.east);
\node[font=\tiny, gray] at (2.0, 1.3) {residual};

% Residual 2: después de ln1 → add2
\draw[skip] (ln1.east) -- ++(1.4,0) -- ++(0,2.3) -- (add2.east);

% ── Anotaciones laterales ─────────────────────────────────────
\node[draw=blue!20, fill=blue!4, rounded corners, font=\tiny,
      inner sep=4pt, align=left] at (-3.2, 1.4)
    {Cada token consulta\\a todos los dem\'as.\\Produce representaci\'on\\contextualizada.};
\draw[-{Stealth}, gray!50, thin] (-1.8, 1.4) -- (-2.1, 1.4);

\node[draw=green!20, fill=green!4, rounded corners, font=\tiny,
      inner sep=4pt, align=left] at (-3.2, 4.6)
    {Transf. no lineal\\aplicada a cada token\\independientemente.\\ReLU o GELU.};
\draw[-{Stealth}, gray!50, thin] (-1.8, 4.6) -- (-2.1, 4.6);

\node[draw=orange!20, fill=orange!4, rounded corners, font=\tiny,
      inner sep=4pt, align=left] at (3.2, 3.4)
    {Normaliza activaciones.\\Estabiliza entrenamiento.};
\draw[-{Stealth}, gray!50, thin] (1.8, 3.4) -- (2.1, 3.4);

% ── Corchete del bloque ───────────────────────────────────────
\draw[gray!50, thick]
    (-2.6, 1.0) -- (-2.8, 1.0) -- (-2.8, 7.0) -- (-2.6, 7.0);
\node[font=\footnotesize, gray, rotate=90] at (-3.1, 4.0)
    {$\times N$ bloques};

\end{tikzpicture}
\caption{Estructura de un bloque encoder. Cada bloque contiene dos
sub-capas: atenci\'on multi-cabeza y una red feed-forward. Ambas est\'an
rodeadas por una conexi\'on residual (\textit{skip connection}) y una
normalizaci\'on de capa (\textit{Layer Norm}). El bloque se apila $N$ veces
para aumentar la capacidad representacional del modelo~\cite{vaswani2017attention}.}
\label{fig:encoder_block}
\end{figure}

% ============================================================
\begin{figure}[H]
\centering
\begin{tikzpicture}[
    block/.style={draw, rounded corners, fill=blue!8, draw=blue!40,
                  minimum width=4.2cm, minimum height=0.7cm,
                  font=\small, align=center},
    norm/.style={draw, rounded corners, fill=orange!10, draw=orange!50,
                 minimum width=4.2cm, minimum height=0.55cm,
                 font=\small, align=center},
    add/.style={draw, circle, fill=yellow!15, draw=yellow!60!black,
                minimum size=0.5cm, font=\small\bfseries},
    arrow/.style={-{Stealth}, thick},
    skip/.style={-{Stealth}, thick, gray!60, dashed}
]

% ── Entrada ──────────────────────────────────────────────────
\node[font=\footnotesize, gray] (inp) at (0, 0) {$\bm{X}$ (embeddings + PE)};

% ── Sub-bloque 1: Multi-Head Self-Attention ──────────────────
\node[block] (mha)  at (0, 1.4) {Multi-Head\\Self-Attention};
\node[add]   (add1) at (0, 2.6) {$+$};
\node[norm]  (ln1)  at (0, 3.4) {Layer Norm};

% ── Sub-bloque 2: Feed-Forward ───────────────────────────────
\node[block] (ffn)  at (0, 4.6) {Feed-Forward\\Network};
\node[add]   (add2) at (0, 5.8) {$+$};
\node[norm]  (ln2)  at (0, 6.6) {Layer Norm};

% ── Salida ───────────────────────────────────────────────────
\node[font=\footnotesize, gray] (out) at (0, 7.6) {$\bm{Z}$ (representaciones contextualizadas)};

% ── Flechas principales ──────────────────────────────────────
\draw[arrow] (inp)  -- (mha);
\draw[arrow] (mha)  -- (add1);
\draw[arrow] (add1) -- (ln1);
\draw[arrow] (ln1)  -- (ffn);
\draw[arrow] (ffn)  -- (add2);
\draw[arrow] (add2) -- (ln2);
\draw[arrow] (ln2)  -- (out);

% ── Conexiones residuales ─────────────────────────────────────
% Residual 1: entrada → add1
\draw[skip] (inp.east) -- ++(1.6,0) -- ++(0,2.6) -- (add1.east);
\node[font=\tiny, gray] at (2.0, 1.3) {residual};

% Residual 2: después de ln1 → add2
\draw[skip] (ln1.east) -- ++(1.4,0) -- ++(0,2.3) -- (add2.east);

% ── Anotaciones laterales ─────────────────────────────────────
\node[draw=blue!20, fill=blue!4, rounded corners, font=\tiny,
      inner sep=4pt, align=left] at (-3.2, 1.4)
    {Cada token consulta\\a todos los dem\'as.\\Produce representaci\'on\\contextualizada.};
\draw[-{Stealth}, gray!50, thin] (-1.8, 1.4) -- (-2.1, 1.4);

\node[draw=green!20, fill=green!4, rounded corners, font=\tiny,
      inner sep=4pt, align=left] at (-3.2, 4.6)
    {Transf. no lineal\\aplicada a cada token\\independientemente.\\ReLU o GELU.};
\draw[-{Stealth}, gray!50, thin] (-1.8, 4.6) -- (-2.1, 4.6);

\node[draw=orange!20, fill=orange!4, rounded corners, font=\tiny,
      inner sep=4pt, align=left] at (3.2, 3.4)
    {Normaliza activaciones.\\Estabiliza entrenamiento.};
\draw[-{Stealth}, gray!50, thin] (1.8, 3.4) -- (2.1, 3.4);

% ── Corchete del bloque ───────────────────────────────────────
\draw[gray!50, thick]
    (-2.6, 1.0) -- (-2.8, 1.0) -- (-2.8, 7.0) -- (-2.6, 7.0);
\node[font=\footnotesize, gray, rotate=90] at (-3.1, 4.0)
    {$\times N$ bloques};

\end{tikzpicture}
\caption{Estructura de un bloque encoder. Cada bloque contiene dos
sub-capas: atenci\'on multi-cabeza y una red feed-forward. Ambas est\'an
rodeadas por una conexi\'on residual (\textit{skip connection}) y una
normalizaci\'on de capa (\textit{Layer Norm}). El bloque se apila $N$ veces
para aumentar la capacidad representacional del modelo~\cite{vaswani2017attention}.}
\label{fig:encoder_block}
\end{figure}

% ============================================================
\begin{figure}[H]
\centering
\begin{tikzpicture}[
    block/.style={draw, rounded corners, fill=blue!8, draw=blue!40,
                  minimum width=4.2cm, minimum height=0.7cm,
                  font=\small, align=center},
    norm/.style={draw, rounded corners, fill=orange!10, draw=orange!50,
                 minimum width=4.2cm, minimum height=0.55cm,
                 font=\small, align=center},
    add/.style={draw, circle, fill=yellow!15, draw=yellow!60!black,
                minimum size=0.5cm, font=\small\bfseries},
    arrow/.style={-{Stealth}, thick},
    skip/.style={-{Stealth}, thick, gray!60, dashed}
]

% ── Entrada ──────────────────────────────────────────────────
\node[font=\footnotesize, gray] (inp) at (0, 0) {$\bm{X}$ (embeddings + PE)};

% ── Sub-bloque 1: Multi-Head Self-Attention ──────────────────
\node[block] (mha)  at (0, 1.4) {Multi-Head\\Self-Attention};
\node[add]   (add1) at (0, 2.6) {$+$};
\node[norm]  (ln1)  at (0, 3.4) {Layer Norm};

% ── Sub-bloque 2: Feed-Forward ───────────────────────────────
\node[block] (ffn)  at (0, 4.6) {Feed-Forward\\Network};
\node[add]   (add2) at (0, 5.8) {$+$};
\node[norm]  (ln2)  at (0, 6.6) {Layer Norm};

% ── Salida ───────────────────────────────────────────────────
\node[font=\footnotesize, gray] (out) at (0, 7.6) {$\bm{Z}$ (representaciones contextualizadas)};

% ── Flechas principales ──────────────────────────────────────
\draw[arrow] (inp)  -- (mha);
\draw[arrow] (mha)  -- (add1);
\draw[arrow] (add1) -- (ln1);
\draw[arrow] (ln1)  -- (ffn);
\draw[arrow] (ffn)  -- (add2);
\draw[arrow] (add2) -- (ln2);
\draw[arrow] (ln2)  -- (out);

% ── Conexiones residuales ─────────────────────────────────────
% Residual 1: entrada → add1
\draw[skip] (inp.east) -- ++(1.6,0) -- ++(0,2.6) -- (add1.east);
\node[font=\tiny, gray] at (2.0, 1.3) {residual};

% Residual 2: después de ln1 → add2
\draw[skip] (ln1.east) -- ++(1.4,0) -- ++(0,2.3) -- (add2.east);

% ── Anotaciones laterales ─────────────────────────────────────
\node[draw=blue!20, fill=blue!4, rounded corners, font=\tiny,
      inner sep=4pt, align=left] at (-3.2, 1.4)
    {Cada token consulta\\a todos los dem\'as.\\Produce representaci\'on\\contextualizada.};
\draw[-{Stealth}, gray!50, thin] (-1.8, 1.4) -- (-2.1, 1.4);

\node[draw=green!20, fill=green!4, rounded corners, font=\tiny,
      inner sep=4pt, align=left] at (-3.2, 4.6)
    {Transf. no lineal\\aplicada a cada token\\independientemente.\\ReLU o GELU.};
\draw[-{Stealth}, gray!50, thin] (-1.8, 4.6) -- (-2.1, 4.6);

\node[draw=orange!20, fill=orange!4, rounded corners, font=\tiny,
      inner sep=4pt, align=left] at (3.2, 3.4)
    {Normaliza activaciones.\\Estabiliza entrenamiento.};
\draw[-{Stealth}, gray!50, thin] (1.8, 3.4) -- (2.1, 3.4);

% ── Corchete del bloque ───────────────────────────────────────
\draw[gray!50, thick]
    (-2.6, 1.0) -- (-2.8, 1.0) -- (-2.8, 7.0) -- (-2.6, 7.0);
\node[font=\footnotesize, gray, rotate=90] at (-3.1, 4.0)
    {$\times N$ bloques};

\end{tikzpicture}
\caption{Estructura de un bloque encoder. Cada bloque contiene dos
sub-capas: atenci\'on multi-cabeza y una red feed-forward. Ambas est\'an
rodeadas por una conexi\'on residual (\textit{skip connection}) y una
normalizaci\'on de capa (\textit{Layer Norm}). El bloque se apila $N$ veces
para aumentar la capacidad representacional del modelo~\cite{vaswani2017attention}.}
\label{fig:encoder_block}
\end{figure}


El encoder completo apila $N$ de estos bloques. Cada bloque puede
especializarse en distintos tipos de relaciones: los primeros bloques
tienden a capturar patrones locales y sintácticos, mientras que los
bloques más profundos capturan relaciones semánticas de mayor
abstracción~\cite{vaswani2017attention}.

%------------------------------------------------------------
\subsection{El bloque decoder}
\label{subsec:bloque_decoder}

El decoder extiende el encoder con un mecanismo adicional para incorporar
información de la secuencia de entrada. Se utiliza en tareas donde la
salida es una secuencia distinta de la entrada, como la traducción
automática o la generación de resúmenes. Consta de tres sub-capas:

\begin{enumerate}
    \item \textbf{Masked Multi-Head Self-Attention:} idéntica a la del
    encoder, pero con una máscara causal que impide que el token en
    posición $t$ acceda a tokens en posiciones $t+1, t+2, \ldots$
    Esto garantiza la propiedad autoregresiva: el modelo genera cada
    token condicionado únicamente en los anteriores.

    \item \textbf{Multi-Head Cross-Attention:} en esta sub-capa, las
    \textit{queries} provienen del decoder (la secuencia generada hasta
    el momento) mientras que las \textit{keys} y \textit{values} provienen
    de la salida del encoder (la representación de la entrada). Esto
    permite al decoder ``consultar'' la secuencia de entrada para decidir
    qué información es relevante en cada paso de generación.

    \item \textbf{Feed-Forward Network:} idéntica a la del encoder.
\end{enumerate}

Las tres sub-capas también están envueltas en conexiones residuales y
Layer Normalization.

%------------------------------------------------------------
\subsection{Arquitectura completa}
\label{subsec:arquitectura_completa}

La Figura~\ref{fig:transformer_completo} muestra la arquitectura completa
del Transformer original, con encoder y decoder.

% ============================================================
% fig_transformer_completo.tex — Arquitectura Transformer completa
% Uso: % ============================================================
% fig_transformer_completo.tex — Arquitectura Transformer completa
% Uso: % ============================================================
% fig_transformer_completo.tex — Arquitectura Transformer completa
% Uso: \input{figures/fig_transformer_completo}
% ============================================================
\begin{figure}[H]
\centering
\begin{tikzpicture}[
    blk/.style={draw, rounded corners, minimum width=3.0cm,
                minimum height=0.62cm, font=\footnotesize, align=center},
    enc/.style={blk, fill=blue!10, draw=blue!40},
    dec/.style={blk, fill=green!10, draw=green!50!black},
    norm/.style={blk, fill=orange!10, draw=orange!50,
                 minimum height=0.5cm},
    emb/.style={blk, fill=purple!10, draw=purple!40},
    out/.style={blk, fill=red!8, draw=red!40},
    add/.style={draw, circle, fill=yellow!20, draw=yellow!60!black,
                minimum size=0.42cm, font=\scriptsize\bfseries},
    arrow/.style={-{Stealth}, thick},
    skip/.style={-{Stealth}, thick, gray!50, dashed},
    lbl/.style={font=\tiny, gray}
]

%═══════════════════════════════════════════════════════════
% ENCODER (izquierda)
%═══════════════════════════════════════════════════════════
\node[font=\small\bfseries, blue!60] at (-3.2, 10.8) {ENCODER};

% Entrada
\node[lbl] (src) at (-3.2, -0.4) {secuencia de entrada};
\node[emb] (src_emb) at (-3.2, 0.55) {Input Embedding};
\node[emb] (src_pe)  at (-3.2, 1.45) {Positional Encoding};

% Bloque encoder (representado como caja grande)
\draw[blue!30, rounded corners, thick, fill=blue!3]
    (-4.9, 2.1) rectangle (-1.5, 7.9);
\node[lbl] at (-4.55, 7.65) {$\times N$};

\node[enc]  (mha_e)  at (-3.2, 2.75) {Multi-Head\\Self-Attention};
\node[add]  (add1e)  at (-3.2, 3.75) {$+$};
\node[norm] (ln1e)   at (-3.2, 4.4)  {Add \& Norm};
\node[enc]  (ffn_e)  at (-3.2, 5.3)  {Feed-Forward};
\node[add]  (add2e)  at (-3.2, 6.3)  {$+$};
\node[norm] (ln2e)   at (-3.2, 6.95) {Add \& Norm};

% Salida encoder
\node[lbl] (enc_out) at (-3.2, 8.3) {$K, V$ al decoder};

% Flechas encoder
\draw[arrow] (src)     -- (src_emb);
\draw[arrow] (src_emb) -- (src_pe);
\draw[arrow] (src_pe)  -- (mha_e);
\draw[arrow] (mha_e)   -- (add1e);
\draw[arrow] (add1e)   -- (ln1e);
\draw[arrow] (ln1e)    -- (ffn_e);
\draw[arrow] (ffn_e)   -- (add2e);
\draw[arrow] (add2e)   -- (ln2e);
\draw[arrow] (ln2e)    -- (enc_out);

% Residuales encoder
\draw[skip] (src_pe.west) -- ++(-0.7,0) -- ++(0,3.75) -- (add1e.west);
\draw[skip] (ln1e.west)   -- ++(-0.5,0) -- ++(0,2.0)  -- (add2e.west);

%═══════════════════════════════════════════════════════════
% DECODER (derecha)
%═══════════════════════════════════════════════════════════
\node[font=\small\bfseries, green!60!black] at (3.2, 10.8) {DECODER};

% Entrada
\node[lbl] (tgt) at (3.2, -0.4) {secuencia de salida (desplazada)};
\node[emb] (tgt_emb) at (3.2, 0.55) {Output Embedding};
\node[emb] (tgt_pe)  at (3.2, 1.45) {Positional Encoding};

% Bloque decoder
\draw[green!30, rounded corners, thick, fill=green!3]
    (1.5, 2.1) rectangle (4.9, 9.1);
\node[lbl] at (4.55, 8.85) {$\times N$};

\node[dec]  (mha_d1) at (3.2, 2.75) {Masked Multi-Head\\Self-Attention};
\node[add]  (add1d)  at (3.2, 3.75) {$+$};
\node[norm] (ln1d)   at (3.2, 4.4)  {Add \& Norm};
\node[dec]  (mha_d2) at (3.2, 5.3)  {Multi-Head\\Cross-Attention};
\node[add]  (add2d)  at (3.2, 6.3)  {$+$};
\node[norm] (ln2d)   at (3.2, 6.95) {Add \& Norm};
\node[dec]  (ffn_d)  at (3.2, 7.85) {Feed-Forward};
\node[add]  (add3d)  at (3.2, 8.75) {$+$};
% (ln3 dentro del bloque, salida)

% Flechas decoder
\draw[arrow] (tgt)     -- (tgt_emb);
\draw[arrow] (tgt_emb) -- (tgt_pe);
\draw[arrow] (tgt_pe)  -- (mha_d1);
\draw[arrow] (mha_d1)  -- (add1d);
\draw[arrow] (add1d)   -- (ln1d);
\draw[arrow] (ln1d)    -- (mha_d2);
\draw[arrow] (mha_d2)  -- (add2d);
\draw[arrow] (add2d)   -- (ln2d);
\draw[arrow] (ln2d)    -- (ffn_d);
\draw[arrow] (ffn_d)   -- (add3d);

% Residuales decoder
\draw[skip] (tgt_pe.east) -- ++(0.7,0) -- ++(0,3.75) -- (add1d.east);
\draw[skip] (ln1d.east)   -- ++(0.5,0) -- ++(0,2.0)  -- (add2d.east);
\draw[skip] (ln2d.east)   -- ++(0.5,0) -- ++(0,2.0)  -- (add3d.east);

% ── Conexión encoder → decoder (cross-attention) ─────────────
\draw[arrow, red!60, very thick]
    (enc_out.east) -- ++(0.5,0) -- ++(0,-3.0) -- (mha_d2.west);
\node[font=\tiny, red!60, align=center] at (0, 6.8)
    {$K$, $V$ del\\encoder};

% ── Salida final ──────────────────────────────────────────────
\node[out]  (linear)  at (3.2, 9.8)  {Linear};
\node[out]  (softmax) at (3.2, 10.5) {Softmax};
\node[lbl]           at (3.2, 11.1)  {probabilidades de salida};

\draw[arrow] (add3d)  -- (linear);
\draw[arrow] (linear)  -- (softmax);

% ── Nota máscara ─────────────────────────────────────────────
\node[draw=gray!30, fill=gray!5, rounded corners, font=\tiny,
      inner sep=4pt, align=left] at (6.8, 2.75)
    {\textbf{Masked:} m\'ascara causal\\impide ver tokens futuros.\\Garantiza generaci\'on\\autoregresiva.};
\draw[-{Stealth}, gray!40, thin] (4.9, 2.75) -- (5.5, 2.75);

\end{tikzpicture}
\caption{Arquitectura completa del Transformer original~\cite{vaswani2017attention}.
El encoder (izquierda) procesa la secuencia de entrada en paralelo y produce
representaciones contextualizadas. El decoder (derecha) genera la secuencia
de salida token a token, utilizando atenci\'on enmascarada sobre su propia
salida parcial y atenci\'on cruzada (\textit{cross-attention}) sobre las
representaciones del encoder. Las flechas rojas indican el flujo de $K$ y $V$
del encoder al decoder.}
\label{fig:transformer_completo}
\end{figure}

% ============================================================
\begin{figure}[H]
\centering
\begin{tikzpicture}[
    blk/.style={draw, rounded corners, minimum width=3.0cm,
                minimum height=0.62cm, font=\footnotesize, align=center},
    enc/.style={blk, fill=blue!10, draw=blue!40},
    dec/.style={blk, fill=green!10, draw=green!50!black},
    norm/.style={blk, fill=orange!10, draw=orange!50,
                 minimum height=0.5cm},
    emb/.style={blk, fill=purple!10, draw=purple!40},
    out/.style={blk, fill=red!8, draw=red!40},
    add/.style={draw, circle, fill=yellow!20, draw=yellow!60!black,
                minimum size=0.42cm, font=\scriptsize\bfseries},
    arrow/.style={-{Stealth}, thick},
    skip/.style={-{Stealth}, thick, gray!50, dashed},
    lbl/.style={font=\tiny, gray}
]

%═══════════════════════════════════════════════════════════
% ENCODER (izquierda)
%═══════════════════════════════════════════════════════════
\node[font=\small\bfseries, blue!60] at (-3.2, 10.8) {ENCODER};

% Entrada
\node[lbl] (src) at (-3.2, -0.4) {secuencia de entrada};
\node[emb] (src_emb) at (-3.2, 0.55) {Input Embedding};
\node[emb] (src_pe)  at (-3.2, 1.45) {Positional Encoding};

% Bloque encoder (representado como caja grande)
\draw[blue!30, rounded corners, thick, fill=blue!3]
    (-4.9, 2.1) rectangle (-1.5, 7.9);
\node[lbl] at (-4.55, 7.65) {$\times N$};

\node[enc]  (mha_e)  at (-3.2, 2.75) {Multi-Head\\Self-Attention};
\node[add]  (add1e)  at (-3.2, 3.75) {$+$};
\node[norm] (ln1e)   at (-3.2, 4.4)  {Add \& Norm};
\node[enc]  (ffn_e)  at (-3.2, 5.3)  {Feed-Forward};
\node[add]  (add2e)  at (-3.2, 6.3)  {$+$};
\node[norm] (ln2e)   at (-3.2, 6.95) {Add \& Norm};

% Salida encoder
\node[lbl] (enc_out) at (-3.2, 8.3) {$K, V$ al decoder};

% Flechas encoder
\draw[arrow] (src)     -- (src_emb);
\draw[arrow] (src_emb) -- (src_pe);
\draw[arrow] (src_pe)  -- (mha_e);
\draw[arrow] (mha_e)   -- (add1e);
\draw[arrow] (add1e)   -- (ln1e);
\draw[arrow] (ln1e)    -- (ffn_e);
\draw[arrow] (ffn_e)   -- (add2e);
\draw[arrow] (add2e)   -- (ln2e);
\draw[arrow] (ln2e)    -- (enc_out);

% Residuales encoder
\draw[skip] (src_pe.west) -- ++(-0.7,0) -- ++(0,3.75) -- (add1e.west);
\draw[skip] (ln1e.west)   -- ++(-0.5,0) -- ++(0,2.0)  -- (add2e.west);

%═══════════════════════════════════════════════════════════
% DECODER (derecha)
%═══════════════════════════════════════════════════════════
\node[font=\small\bfseries, green!60!black] at (3.2, 10.8) {DECODER};

% Entrada
\node[lbl] (tgt) at (3.2, -0.4) {secuencia de salida (desplazada)};
\node[emb] (tgt_emb) at (3.2, 0.55) {Output Embedding};
\node[emb] (tgt_pe)  at (3.2, 1.45) {Positional Encoding};

% Bloque decoder
\draw[green!30, rounded corners, thick, fill=green!3]
    (1.5, 2.1) rectangle (4.9, 9.1);
\node[lbl] at (4.55, 8.85) {$\times N$};

\node[dec]  (mha_d1) at (3.2, 2.75) {Masked Multi-Head\\Self-Attention};
\node[add]  (add1d)  at (3.2, 3.75) {$+$};
\node[norm] (ln1d)   at (3.2, 4.4)  {Add \& Norm};
\node[dec]  (mha_d2) at (3.2, 5.3)  {Multi-Head\\Cross-Attention};
\node[add]  (add2d)  at (3.2, 6.3)  {$+$};
\node[norm] (ln2d)   at (3.2, 6.95) {Add \& Norm};
\node[dec]  (ffn_d)  at (3.2, 7.85) {Feed-Forward};
\node[add]  (add3d)  at (3.2, 8.75) {$+$};
% (ln3 dentro del bloque, salida)

% Flechas decoder
\draw[arrow] (tgt)     -- (tgt_emb);
\draw[arrow] (tgt_emb) -- (tgt_pe);
\draw[arrow] (tgt_pe)  -- (mha_d1);
\draw[arrow] (mha_d1)  -- (add1d);
\draw[arrow] (add1d)   -- (ln1d);
\draw[arrow] (ln1d)    -- (mha_d2);
\draw[arrow] (mha_d2)  -- (add2d);
\draw[arrow] (add2d)   -- (ln2d);
\draw[arrow] (ln2d)    -- (ffn_d);
\draw[arrow] (ffn_d)   -- (add3d);

% Residuales decoder
\draw[skip] (tgt_pe.east) -- ++(0.7,0) -- ++(0,3.75) -- (add1d.east);
\draw[skip] (ln1d.east)   -- ++(0.5,0) -- ++(0,2.0)  -- (add2d.east);
\draw[skip] (ln2d.east)   -- ++(0.5,0) -- ++(0,2.0)  -- (add3d.east);

% ── Conexión encoder → decoder (cross-attention) ─────────────
\draw[arrow, red!60, very thick]
    (enc_out.east) -- ++(0.5,0) -- ++(0,-3.0) -- (mha_d2.west);
\node[font=\tiny, red!60, align=center] at (0, 6.8)
    {$K$, $V$ del\\encoder};

% ── Salida final ──────────────────────────────────────────────
\node[out]  (linear)  at (3.2, 9.8)  {Linear};
\node[out]  (softmax) at (3.2, 10.5) {Softmax};
\node[lbl]           at (3.2, 11.1)  {probabilidades de salida};

\draw[arrow] (add3d)  -- (linear);
\draw[arrow] (linear)  -- (softmax);

% ── Nota máscara ─────────────────────────────────────────────
\node[draw=gray!30, fill=gray!5, rounded corners, font=\tiny,
      inner sep=4pt, align=left] at (6.8, 2.75)
    {\textbf{Masked:} m\'ascara causal\\impide ver tokens futuros.\\Garantiza generaci\'on\\autoregresiva.};
\draw[-{Stealth}, gray!40, thin] (4.9, 2.75) -- (5.5, 2.75);

\end{tikzpicture}
\caption{Arquitectura completa del Transformer original~\cite{vaswani2017attention}.
El encoder (izquierda) procesa la secuencia de entrada en paralelo y produce
representaciones contextualizadas. El decoder (derecha) genera la secuencia
de salida token a token, utilizando atenci\'on enmascarada sobre su propia
salida parcial y atenci\'on cruzada (\textit{cross-attention}) sobre las
representaciones del encoder. Las flechas rojas indican el flujo de $K$ y $V$
del encoder al decoder.}
\label{fig:transformer_completo}
\end{figure}

% ============================================================
\begin{figure}[H]
\centering
\begin{tikzpicture}[
    blk/.style={draw, rounded corners, minimum width=3.0cm,
                minimum height=0.62cm, font=\footnotesize, align=center},
    enc/.style={blk, fill=blue!10, draw=blue!40},
    dec/.style={blk, fill=green!10, draw=green!50!black},
    norm/.style={blk, fill=orange!10, draw=orange!50,
                 minimum height=0.5cm},
    emb/.style={blk, fill=purple!10, draw=purple!40},
    out/.style={blk, fill=red!8, draw=red!40},
    add/.style={draw, circle, fill=yellow!20, draw=yellow!60!black,
                minimum size=0.42cm, font=\scriptsize\bfseries},
    arrow/.style={-{Stealth}, thick},
    skip/.style={-{Stealth}, thick, gray!50, dashed},
    lbl/.style={font=\tiny, gray}
]

%═══════════════════════════════════════════════════════════
% ENCODER (izquierda)
%═══════════════════════════════════════════════════════════
\node[font=\small\bfseries, blue!60] at (-3.2, 10.8) {ENCODER};

% Entrada
\node[lbl] (src) at (-3.2, -0.4) {secuencia de entrada};
\node[emb] (src_emb) at (-3.2, 0.55) {Input Embedding};
\node[emb] (src_pe)  at (-3.2, 1.45) {Positional Encoding};

% Bloque encoder (representado como caja grande)
\draw[blue!30, rounded corners, thick, fill=blue!3]
    (-4.9, 2.1) rectangle (-1.5, 7.9);
\node[lbl] at (-4.55, 7.65) {$\times N$};

\node[enc]  (mha_e)  at (-3.2, 2.75) {Multi-Head\\Self-Attention};
\node[add]  (add1e)  at (-3.2, 3.75) {$+$};
\node[norm] (ln1e)   at (-3.2, 4.4)  {Add \& Norm};
\node[enc]  (ffn_e)  at (-3.2, 5.3)  {Feed-Forward};
\node[add]  (add2e)  at (-3.2, 6.3)  {$+$};
\node[norm] (ln2e)   at (-3.2, 6.95) {Add \& Norm};

% Salida encoder
\node[lbl] (enc_out) at (-3.2, 8.3) {$K, V$ al decoder};

% Flechas encoder
\draw[arrow] (src)     -- (src_emb);
\draw[arrow] (src_emb) -- (src_pe);
\draw[arrow] (src_pe)  -- (mha_e);
\draw[arrow] (mha_e)   -- (add1e);
\draw[arrow] (add1e)   -- (ln1e);
\draw[arrow] (ln1e)    -- (ffn_e);
\draw[arrow] (ffn_e)   -- (add2e);
\draw[arrow] (add2e)   -- (ln2e);
\draw[arrow] (ln2e)    -- (enc_out);

% Residuales encoder
\draw[skip] (src_pe.west) -- ++(-0.7,0) -- ++(0,3.75) -- (add1e.west);
\draw[skip] (ln1e.west)   -- ++(-0.5,0) -- ++(0,2.0)  -- (add2e.west);

%═══════════════════════════════════════════════════════════
% DECODER (derecha)
%═══════════════════════════════════════════════════════════
\node[font=\small\bfseries, green!60!black] at (3.2, 10.8) {DECODER};

% Entrada
\node[lbl] (tgt) at (3.2, -0.4) {secuencia de salida (desplazada)};
\node[emb] (tgt_emb) at (3.2, 0.55) {Output Embedding};
\node[emb] (tgt_pe)  at (3.2, 1.45) {Positional Encoding};

% Bloque decoder
\draw[green!30, rounded corners, thick, fill=green!3]
    (1.5, 2.1) rectangle (4.9, 9.1);
\node[lbl] at (4.55, 8.85) {$\times N$};

\node[dec]  (mha_d1) at (3.2, 2.75) {Masked Multi-Head\\Self-Attention};
\node[add]  (add1d)  at (3.2, 3.75) {$+$};
\node[norm] (ln1d)   at (3.2, 4.4)  {Add \& Norm};
\node[dec]  (mha_d2) at (3.2, 5.3)  {Multi-Head\\Cross-Attention};
\node[add]  (add2d)  at (3.2, 6.3)  {$+$};
\node[norm] (ln2d)   at (3.2, 6.95) {Add \& Norm};
\node[dec]  (ffn_d)  at (3.2, 7.85) {Feed-Forward};
\node[add]  (add3d)  at (3.2, 8.75) {$+$};
% (ln3 dentro del bloque, salida)

% Flechas decoder
\draw[arrow] (tgt)     -- (tgt_emb);
\draw[arrow] (tgt_emb) -- (tgt_pe);
\draw[arrow] (tgt_pe)  -- (mha_d1);
\draw[arrow] (mha_d1)  -- (add1d);
\draw[arrow] (add1d)   -- (ln1d);
\draw[arrow] (ln1d)    -- (mha_d2);
\draw[arrow] (mha_d2)  -- (add2d);
\draw[arrow] (add2d)   -- (ln2d);
\draw[arrow] (ln2d)    -- (ffn_d);
\draw[arrow] (ffn_d)   -- (add3d);

% Residuales decoder
\draw[skip] (tgt_pe.east) -- ++(0.7,0) -- ++(0,3.75) -- (add1d.east);
\draw[skip] (ln1d.east)   -- ++(0.5,0) -- ++(0,2.0)  -- (add2d.east);
\draw[skip] (ln2d.east)   -- ++(0.5,0) -- ++(0,2.0)  -- (add3d.east);

% ── Conexión encoder → decoder (cross-attention) ─────────────
\draw[arrow, red!60, very thick]
    (enc_out.east) -- ++(0.5,0) -- ++(0,-3.0) -- (mha_d2.west);
\node[font=\tiny, red!60, align=center] at (0, 6.8)
    {$K$, $V$ del\\encoder};

% ── Salida final ──────────────────────────────────────────────
\node[out]  (linear)  at (3.2, 9.8)  {Linear};
\node[out]  (softmax) at (3.2, 10.5) {Softmax};
\node[lbl]           at (3.2, 11.1)  {probabilidades de salida};

\draw[arrow] (add3d)  -- (linear);
\draw[arrow] (linear)  -- (softmax);

% ── Nota máscara ─────────────────────────────────────────────
\node[draw=gray!30, fill=gray!5, rounded corners, font=\tiny,
      inner sep=4pt, align=left] at (6.8, 2.75)
    {\textbf{Masked:} m\'ascara causal\\impide ver tokens futuros.\\Garantiza generaci\'on\\autoregresiva.};
\draw[-{Stealth}, gray!40, thin] (4.9, 2.75) -- (5.5, 2.75);

\end{tikzpicture}
\caption{Arquitectura completa del Transformer original~\cite{vaswani2017attention}.
El encoder (izquierda) procesa la secuencia de entrada en paralelo y produce
representaciones contextualizadas. El decoder (derecha) genera la secuencia
de salida token a token, utilizando atenci\'on enmascarada sobre su propia
salida parcial y atenci\'on cruzada (\textit{cross-attention}) sobre las
representaciones del encoder. Las flechas rojas indican el flujo de $K$ y $V$
del encoder al decoder.}
\label{fig:transformer_completo}
\end{figure}


El flujo de información puede resumirse así: la secuencia de entrada pasa
por el encoder y produce un conjunto de representaciones ricas en contexto.
El decoder recibe estas representaciones a través del cross-attention y,
combinándolas con su propio estado interno, genera la secuencia de salida
token a token. En cada paso, el decoder produce una distribución de
probabilidad sobre el vocabulario y selecciona el token más probable (o
muestrea de la distribución), que luego se alimenta como entrada al
siguiente paso.

%------------------------------------------------------------
\subsection{Variantes principales}
\label{subsec:variantes}

La arquitectura original de Vaswani et al.\ incluye encoder y decoder.
Sin embargo, muchas aplicaciones prácticas utilizan únicamente uno de
los dos componentes, dependiendo de la tarea:

% ============================================================
% fig_variantes_transformer.tex — Variantes encoder-only vs enc-dec
% Uso: % ============================================================
% fig_variantes_transformer.tex — Variantes encoder-only vs enc-dec
% Uso: % ============================================================
% fig_variantes_transformer.tex — Variantes encoder-only vs enc-dec
% Uso: \input{figures/fig_variantes_transformer}
% ============================================================
\begin{figure}[H]
\centering
\begin{tikzpicture}[
    hdr/.style={font=\small\bfseries, align=center},
    cell/.style={font=\footnotesize, align=center},
    enconly/.style={draw=blue!40, fill=blue!8, rounded corners,
                    minimum width=2.2cm, minimum height=0.8cm,
                    font=\footnotesize\bfseries, align=center},
    encdec/.style={draw=green!50!black, fill=green!8, rounded corners,
                   minimum width=2.2cm, minimum height=0.8cm,
                   font=\footnotesize\bfseries, align=center}
]

% ── Mini diagramas ───────────────────────────────────────────

% Encoder-only
\node[enconly] at (1.5, 6.5) {Encoder\\$\times N$};
\node[font=\tiny, gray] at (1.5, 5.9) {entrada};
\node[font=\tiny, gray] at (1.5, 7.1) {repr. contextual};
\draw[-{Stealth}, thick, blue!50] (1.5, 6.1) -- (1.5, 6.15);
\draw[-{Stealth}, thick, blue!50] (1.5, 6.85) -- (1.5, 6.9);
\node[font=\small\bfseries, blue!60] at (1.5, 7.5) {Encoder-only};
\node[font=\tiny, gray] at (1.5, 7.9) {(ej: BERT)};

% Encoder-Decoder
\node[enconly]  at (6.5, 6.5) {Encoder\\$\times N$};
\node[encdec]   at (8.8, 6.5) {Decoder\\$\times N$};
\draw[-{Stealth}, thick, red!50] (7.7, 6.5) -- (8.05, 6.5);
\node[font=\tiny, red!60] at (7.87, 6.75) {$K,V$};
\node[font=\tiny, gray]   at (6.5, 5.9) {entrada};
\node[font=\tiny, gray]   at (8.8, 7.1) {salida};
\draw[-{Stealth}, thick, blue!50]  (6.5, 6.1) -- (6.5, 6.15);
\draw[-{Stealth}, thick, green!50!black] (8.8, 6.85) -- (8.8, 6.9);
\node[font=\small\bfseries, green!60!black] at (7.65, 7.5) {Encoder-Decoder};
\node[font=\tiny, gray] at (7.65, 7.9) {(ej: T5, BART)};

% Decoder-only
\node[encdec] at (12.5, 6.5) {Decoder\\(causal) $\times N$};
\node[font=\tiny, gray] at (12.5, 5.9) {contexto previo};
\node[font=\tiny, gray] at (12.5, 7.1) {siguiente token};
\draw[-{Stealth}, thick, green!50!black] (12.5, 6.1) -- (12.5, 6.15);
\draw[-{Stealth}, thick, green!50!black] (12.5, 6.85) -- (12.5, 6.9);
\node[font=\small\bfseries, green!60!black] at (12.5, 7.5) {Decoder-only};
\node[font=\tiny, gray] at (12.5, 7.9) {(ej: GPT, LLaMA)};

% ── Tabla comparativa ─────────────────────────────────────────
\node[hdr] at (0,    5.0) {Caracter\'istica};
\node[hdr, blue!60] at (3.5,  5.0) {Encoder-only};
\node[hdr, green!60!black] at (7.5,  5.0) {Encoder-Decoder};
\node[hdr, green!60!black] at (11.5, 5.0) {Decoder-only};

\draw[gray!40, thick] (-0.5, 4.7) -- (13.5, 4.7);

% Filas
\foreach \y/\lbl/\a/\b/\c in {
    4.1/{Direcci\'on atenci\'on}/{bidireccional}/{bidireccional (enc)}/{unidireccional},
    3.4/{Tarea principal}/{clasificaci\'on, NER}/{traducci\'on, res\'umenes}/{generaci\'on de texto},
    2.7/{Ve tokens futuros}/{S\'i}/{S\'i (encoder)}/{No (m\'ascara causal)},
    2.0/{Tama\~no relativo}/{mediano}/{grande}/{grande}
}{
    \node[cell, align=left] at (0, \y) {\lbl};
    \node[cell] at (3.5,  \y) {\a};
    \node[cell] at (7.5,  \y) {\b};
    \node[cell] at (11.5, \y) {\c};
    \draw[gray!20] (-0.5, \y-0.35) -- (13.5, \y-0.35);
}

\draw[gray!40] (-0.5, 5.3) -- (-0.5, 1.65);
\draw[gray!40] (2.0,  5.3) -- (2.0,  1.65);
\draw[gray!40] (5.5,  5.3) -- (5.5,  1.65);
\draw[gray!40] (9.5,  5.3) -- (9.5,  1.65);
\draw[gray!40] (13.5, 5.3) -- (13.5, 1.65);

\end{tikzpicture}
\caption{Las tres variantes principales de la arquitectura Transformer.
\textbf{Encoder-only} (BERT): procesa la entrada completa en ambas
direcciones; ideal para tareas de comprensi\'on. \textbf{Encoder-decoder}
(T5, BART): mapea una secuencia a otra; ideal para traducci\'on y res\'umenes.
\textbf{Decoder-only} (GPT, LLaMA): genera texto token a token con m\'ascara
causal; es la base de los modelos de lenguaje actuales.}
\label{fig:variantes_transformer}
\end{figure}

% ============================================================
\begin{figure}[H]
\centering
\begin{tikzpicture}[
    hdr/.style={font=\small\bfseries, align=center},
    cell/.style={font=\footnotesize, align=center},
    enconly/.style={draw=blue!40, fill=blue!8, rounded corners,
                    minimum width=2.2cm, minimum height=0.8cm,
                    font=\footnotesize\bfseries, align=center},
    encdec/.style={draw=green!50!black, fill=green!8, rounded corners,
                   minimum width=2.2cm, minimum height=0.8cm,
                   font=\footnotesize\bfseries, align=center}
]

% ── Mini diagramas ───────────────────────────────────────────

% Encoder-only
\node[enconly] at (1.5, 6.5) {Encoder\\$\times N$};
\node[font=\tiny, gray] at (1.5, 5.9) {entrada};
\node[font=\tiny, gray] at (1.5, 7.1) {repr. contextual};
\draw[-{Stealth}, thick, blue!50] (1.5, 6.1) -- (1.5, 6.15);
\draw[-{Stealth}, thick, blue!50] (1.5, 6.85) -- (1.5, 6.9);
\node[font=\small\bfseries, blue!60] at (1.5, 7.5) {Encoder-only};
\node[font=\tiny, gray] at (1.5, 7.9) {(ej: BERT)};

% Encoder-Decoder
\node[enconly]  at (6.5, 6.5) {Encoder\\$\times N$};
\node[encdec]   at (8.8, 6.5) {Decoder\\$\times N$};
\draw[-{Stealth}, thick, red!50] (7.7, 6.5) -- (8.05, 6.5);
\node[font=\tiny, red!60] at (7.87, 6.75) {$K,V$};
\node[font=\tiny, gray]   at (6.5, 5.9) {entrada};
\node[font=\tiny, gray]   at (8.8, 7.1) {salida};
\draw[-{Stealth}, thick, blue!50]  (6.5, 6.1) -- (6.5, 6.15);
\draw[-{Stealth}, thick, green!50!black] (8.8, 6.85) -- (8.8, 6.9);
\node[font=\small\bfseries, green!60!black] at (7.65, 7.5) {Encoder-Decoder};
\node[font=\tiny, gray] at (7.65, 7.9) {(ej: T5, BART)};

% Decoder-only
\node[encdec] at (12.5, 6.5) {Decoder\\(causal) $\times N$};
\node[font=\tiny, gray] at (12.5, 5.9) {contexto previo};
\node[font=\tiny, gray] at (12.5, 7.1) {siguiente token};
\draw[-{Stealth}, thick, green!50!black] (12.5, 6.1) -- (12.5, 6.15);
\draw[-{Stealth}, thick, green!50!black] (12.5, 6.85) -- (12.5, 6.9);
\node[font=\small\bfseries, green!60!black] at (12.5, 7.5) {Decoder-only};
\node[font=\tiny, gray] at (12.5, 7.9) {(ej: GPT, LLaMA)};

% ── Tabla comparativa ─────────────────────────────────────────
\node[hdr] at (0,    5.0) {Caracter\'istica};
\node[hdr, blue!60] at (3.5,  5.0) {Encoder-only};
\node[hdr, green!60!black] at (7.5,  5.0) {Encoder-Decoder};
\node[hdr, green!60!black] at (11.5, 5.0) {Decoder-only};

\draw[gray!40, thick] (-0.5, 4.7) -- (13.5, 4.7);

% Filas
\foreach \y/\lbl/\a/\b/\c in {
    4.1/{Direcci\'on atenci\'on}/{bidireccional}/{bidireccional (enc)}/{unidireccional},
    3.4/{Tarea principal}/{clasificaci\'on, NER}/{traducci\'on, res\'umenes}/{generaci\'on de texto},
    2.7/{Ve tokens futuros}/{S\'i}/{S\'i (encoder)}/{No (m\'ascara causal)},
    2.0/{Tama\~no relativo}/{mediano}/{grande}/{grande}
}{
    \node[cell, align=left] at (0, \y) {\lbl};
    \node[cell] at (3.5,  \y) {\a};
    \node[cell] at (7.5,  \y) {\b};
    \node[cell] at (11.5, \y) {\c};
    \draw[gray!20] (-0.5, \y-0.35) -- (13.5, \y-0.35);
}

\draw[gray!40] (-0.5, 5.3) -- (-0.5, 1.65);
\draw[gray!40] (2.0,  5.3) -- (2.0,  1.65);
\draw[gray!40] (5.5,  5.3) -- (5.5,  1.65);
\draw[gray!40] (9.5,  5.3) -- (9.5,  1.65);
\draw[gray!40] (13.5, 5.3) -- (13.5, 1.65);

\end{tikzpicture}
\caption{Las tres variantes principales de la arquitectura Transformer.
\textbf{Encoder-only} (BERT): procesa la entrada completa en ambas
direcciones; ideal para tareas de comprensi\'on. \textbf{Encoder-decoder}
(T5, BART): mapea una secuencia a otra; ideal para traducci\'on y res\'umenes.
\textbf{Decoder-only} (GPT, LLaMA): genera texto token a token con m\'ascara
causal; es la base de los modelos de lenguaje actuales.}
\label{fig:variantes_transformer}
\end{figure}

% ============================================================
\begin{figure}[H]
\centering
\begin{tikzpicture}[
    hdr/.style={font=\small\bfseries, align=center},
    cell/.style={font=\footnotesize, align=center},
    enconly/.style={draw=blue!40, fill=blue!8, rounded corners,
                    minimum width=2.2cm, minimum height=0.8cm,
                    font=\footnotesize\bfseries, align=center},
    encdec/.style={draw=green!50!black, fill=green!8, rounded corners,
                   minimum width=2.2cm, minimum height=0.8cm,
                   font=\footnotesize\bfseries, align=center}
]

% ── Mini diagramas ───────────────────────────────────────────

% Encoder-only
\node[enconly] at (1.5, 6.5) {Encoder\\$\times N$};
\node[font=\tiny, gray] at (1.5, 5.9) {entrada};
\node[font=\tiny, gray] at (1.5, 7.1) {repr. contextual};
\draw[-{Stealth}, thick, blue!50] (1.5, 6.1) -- (1.5, 6.15);
\draw[-{Stealth}, thick, blue!50] (1.5, 6.85) -- (1.5, 6.9);
\node[font=\small\bfseries, blue!60] at (1.5, 7.5) {Encoder-only};
\node[font=\tiny, gray] at (1.5, 7.9) {(ej: BERT)};

% Encoder-Decoder
\node[enconly]  at (6.5, 6.5) {Encoder\\$\times N$};
\node[encdec]   at (8.8, 6.5) {Decoder\\$\times N$};
\draw[-{Stealth}, thick, red!50] (7.7, 6.5) -- (8.05, 6.5);
\node[font=\tiny, red!60] at (7.87, 6.75) {$K,V$};
\node[font=\tiny, gray]   at (6.5, 5.9) {entrada};
\node[font=\tiny, gray]   at (8.8, 7.1) {salida};
\draw[-{Stealth}, thick, blue!50]  (6.5, 6.1) -- (6.5, 6.15);
\draw[-{Stealth}, thick, green!50!black] (8.8, 6.85) -- (8.8, 6.9);
\node[font=\small\bfseries, green!60!black] at (7.65, 7.5) {Encoder-Decoder};
\node[font=\tiny, gray] at (7.65, 7.9) {(ej: T5, BART)};

% Decoder-only
\node[encdec] at (12.5, 6.5) {Decoder\\(causal) $\times N$};
\node[font=\tiny, gray] at (12.5, 5.9) {contexto previo};
\node[font=\tiny, gray] at (12.5, 7.1) {siguiente token};
\draw[-{Stealth}, thick, green!50!black] (12.5, 6.1) -- (12.5, 6.15);
\draw[-{Stealth}, thick, green!50!black] (12.5, 6.85) -- (12.5, 6.9);
\node[font=\small\bfseries, green!60!black] at (12.5, 7.5) {Decoder-only};
\node[font=\tiny, gray] at (12.5, 7.9) {(ej: GPT, LLaMA)};

% ── Tabla comparativa ─────────────────────────────────────────
\node[hdr] at (0,    5.0) {Caracter\'istica};
\node[hdr, blue!60] at (3.5,  5.0) {Encoder-only};
\node[hdr, green!60!black] at (7.5,  5.0) {Encoder-Decoder};
\node[hdr, green!60!black] at (11.5, 5.0) {Decoder-only};

\draw[gray!40, thick] (-0.5, 4.7) -- (13.5, 4.7);

% Filas
\foreach \y/\lbl/\a/\b/\c in {
    4.1/{Direcci\'on atenci\'on}/{bidireccional}/{bidireccional (enc)}/{unidireccional},
    3.4/{Tarea principal}/{clasificaci\'on, NER}/{traducci\'on, res\'umenes}/{generaci\'on de texto},
    2.7/{Ve tokens futuros}/{S\'i}/{S\'i (encoder)}/{No (m\'ascara causal)},
    2.0/{Tama\~no relativo}/{mediano}/{grande}/{grande}
}{
    \node[cell, align=left] at (0, \y) {\lbl};
    \node[cell] at (3.5,  \y) {\a};
    \node[cell] at (7.5,  \y) {\b};
    \node[cell] at (11.5, \y) {\c};
    \draw[gray!20] (-0.5, \y-0.35) -- (13.5, \y-0.35);
}

\draw[gray!40] (-0.5, 5.3) -- (-0.5, 1.65);
\draw[gray!40] (2.0,  5.3) -- (2.0,  1.65);
\draw[gray!40] (5.5,  5.3) -- (5.5,  1.65);
\draw[gray!40] (9.5,  5.3) -- (9.5,  1.65);
\draw[gray!40] (13.5, 5.3) -- (13.5, 1.65);

\end{tikzpicture}
\caption{Las tres variantes principales de la arquitectura Transformer.
\textbf{Encoder-only} (BERT): procesa la entrada completa en ambas
direcciones; ideal para tareas de comprensi\'on. \textbf{Encoder-decoder}
(T5, BART): mapea una secuencia a otra; ideal para traducci\'on y res\'umenes.
\textbf{Decoder-only} (GPT, LLaMA): genera texto token a token con m\'ascara
causal; es la base de los modelos de lenguaje actuales.}
\label{fig:variantes_transformer}
\end{figure}


Las tres variantes comparten los mismos bloques fundamentales —embeddings,
codificación posicional, atención multi-cabeza, FFN, conexiones residuales
y Layer Norm— pero difieren en cómo los ensamblan y en qué dirección
procesan la información. Esta modularidad es una de las razones del éxito
del Transformer como arquitectura universal.

%------------------------------------------------------------
\subsection{Hiperparámetros principales}
\label{subsec:hiperparametros}

El comportamiento del Transformer está controlado por un conjunto de
hiperparámetros que determinan su capacidad y costo computacional. La
Tabla~\ref{tab:hiperparametros} resume los más importantes.

\begin{table}[H]
\centering
\caption{Hiperparámetros principales del Transformer y su efecto.}
\label{tab:hiperparametros}
\begin{tabularx}{\textwidth}{llX}
\toprule
\textbf{Parámetro} & \textbf{Notación} & \textbf{Efecto} \\
\midrule
Dimensión del modelo   & $d_{\text{model}}$ & Tamaño de todos los vectores internos. Afecta capacidad y memoria. \\
Número de cabezas      & $h$               & Cuántos tipos de relaciones puede capturar simultáneamente. \\
Número de bloques      & $N$               & Profundidad del modelo. Más bloques = más abstracción. \\
Dimensión FFN          & $d_{\text{ff}}$   & Generalmente $4 \times d_{\text{model}}$. Controla la ``memoria'' por posición. \\
Longitud máx.\ de secuencia & $n_{\text{max}}$ & Límite de contexto. El costo de atención escala como $O(n^2)$. \\
\bottomrule
\end{tabularx}
\end{table}

Para el prototipo mínimo implementado en este trabajo (descrito en la
Sección~\ref{sec:transformer_minimo}), se utilizan valores reducidos de
estos hiperparámetros que permiten el entrenamiento en hardware doméstico
sin sacrificar la comprensión de los mecanismos.
